% Part: first-order-logic
% Chapter: introduction
% Section: semantic-notions

\documentclass[../../../include/open-logic-section]{subfiles}

\begin{document}

\olfileid[fr]{fol}{int}{sem}

\olsection{Notions sémantiques}

Nous avons indiqué plus haut que, lorsque nous examinons si
$\Sat{M}{!A}[s]$ vaut, nous convenons par commodité que~$s$ attribue
des valeurs à toutes les !!{variable}s, mais que seules les valeurs
attribuées aux !!{variable}s qui figurent dans~$!A$ sont utilisées. En
réalité, seules comptent les valeurs des !!{variable}s \emph{libres}
dans~$!A$. Comme nous veillons à tout justifier, nous démontrerons ce
fait. Puisque les !!{sentence}s n'ont aucune variable libre, le choix
de~$s$ ne joue aucun rôle dans leur satisfaction par
!!a{structure}. Lorsque $!A$ est !!a{sentence}, nous pouvons donc
définir $\Sat{M}{!A}$ comme signifiant que
« $\Sat{M}{!A}[s]$ vaut pour toute assignation~$s$ ». Or cette
condition équivaut à l'existence d'au moins une assignation~$s$ telle
que $\Sat{M}{!A}[s]$. Nous devons introduire les assignations de
!!{variable}s pour obtenir une définition convenable de la satisfaction
des !!{formula}s; pour les !!{sentence}s, en revanche, la satisfaction
est indépendante de ces assignations.

Une fois $\Sat{M}{!A}$ défini, nous savons ce que signifient « cas » et
« vrai dans » pour les !!{sentence}s de la logique du premier ordre.
À partir de la définition de $\Sat{M}{!A}$ pour les !!{sentence}s, nous
pouvons alors définir les notions sémantiques fondamentales de
validité, de conséquence logique et de satisfaisabilité.
!!^a{sentence} est \emph{valide}, $\Entails !A$, si toute
!!{structure} le satisfait. L'!!{sentence}~$!A$ est une
\emph{conséquence logique} d'un ensemble d'!!{sentence}s,
$\Gamma \Entails !A$, si toute !!{structure} qui satisfait tous les
!!{sentence}s de~$\Gamma$ satisfait aussi~$!A$. Enfin, un ensemble
d'!!{sentence}s est \emph{satisfaisable} s'il existe
!!a{structure} qui satisfait simultanément tous ses !!{sentence}s.

Les !!{formula}s étant définies inductivement, et leur satisfaction
étant à son tour définie par induction sur leur structure, nous pouvons
recourir à l'induction pour démontrer des propriétés de notre
sémantique et relier entre elles les notions sémantiques que nous
venons de définir. Nous rassemblerons et démontrerons certaines de ces
propriétés, en partie parce qu'elles sont intéressantes en elles-mêmes,
mais surtout parce que beaucoup nous seront utiles lorsque nous
étudierons la relation entre la sémantique et les systèmes de
!!{derivation}. Pour cela, nous devrons aussi définir avec précision,
c'est-à-dire par induction, plusieurs notions et opérations syntaxiques
que nous n'avons pas encore mentionnées.

\end{document}
